\chapter{被擄歸回總結}

\section{瑪拉基書中的婚姻觀與以色列人婚姻歷史的關係}

\subsection{經文}
\subsubsection{瑪拉基書2:10-16}
我們豈不都是一位父嗎？豈不是一位神所造的嗎？我們各人怎麼以詭詐待弟兄，背棄了神與我們列祖所立的約呢？猶大人行事詭詐，並且在以色列和耶路撒冷中行一件可憎的事；因為猶大人褻瀆耶和華所喜愛的聖潔（或作：聖地）,娶事奉外邦神的女子為妻。凡行這事的，無論何人，就是獻供物給萬軍之耶和華，耶和華也必從雅各的帳棚中剪除他。你們又行了一件這樣的事，使前妻歎息哭泣的眼淚遮蓋耶和華的壇，以致耶和華不再看顧那供物，也不樂意從你們手中收納。 你們還說：這是為什麼呢？」因耶和華在你和你幼年所娶的妻中間作見證。他雖是你的配偶，又是你盟約的妻，你卻以詭詐待他。雖然神有靈的餘力能造多人，他不是單造一人嗎？為何只造一人呢？乃是他願人得虔誠的後裔。所以當謹守你們的心，誰也不可以詭詐待幼年所娶的妻。耶和華以色列的神說：「休妻的事和以強暴待妻的人都是我所恨惡的！所以當謹守你們的心，不可行詭詐。」這是萬軍之耶和華說的。 
\subsubsection{經文的大綱} 
1. 先知責備猶大人在婚姻生活上背棄了神與列祖所立的約 (瑪2：10） 
2. 猶大人婚姻中第一件背約的事  (瑪2：11-12） 
a. 猶大人娶事奉外邦神的女子為妻 
b. 神的態度必從雅各的帳棚中剪除行這事的人 
3. 猶大人婚姻中第二件背約的事  (瑪2：13-16） 
a. 猶大人離棄幼年所娶的妻  (瑪2：13a） 
b. 神的態度 (瑪2：13b - 16） 
1) 不再看顧那供物，也不樂意從 猶大人手中收納  (瑪2：13b） 
2) 神不悅納的原因 (瑪2：14-15） 
3) 神對人的要求： 謹守自己的心，不可行詭詐 (瑪2：16）

\subsection{以色列人的婚姻的歷史觀及神對婚姻的盟約教導}
\subsubsection{從瑪拉基書分析以色列人婚姻的歷史觀念}
\subsubsubsection{婚姻的來源}
從瑪拉基書不難看出以色列人認為，婚姻是神所設立的。他們的婚姻觀源於摩西五經的《
創世記》。“雖然神有靈的餘力能造多人，他不是單造一人嗎？” （瑪2：15 a）。 從這節
經文，可以清楚的看出，瑪拉基在這裡所指的“單造一人”指的就是創世紀中所記載的亞
當。“耶和華神用地上的塵土造人，將生氣吹在他鼻孔裡，他就成了有靈的活人，名叫亞
當”（創2：7）。

神在伊甸園裡創造亞當以後，在人類諸多的直系關係（夫妻關係，父母關係，子女關係，
弟兄姐妹關係）中設立的第一個關係就是婚姻。婚姻的設立過程，記載在《創世紀》的第
二章。“耶和華神說：那人獨居不好，我要為他造一個配偶幫助他。。。只是那人沒有遇
見配偶幫助他。耶和華神使他沉睡，他就睡了；於是取下他的一條肋骨，又把肉合起來。
耶和華神就用那人身上所取的肋骨造成一個女人，領他到那人跟前。那人說：這是我骨中
的骨，肉中的肉，可以稱他為女人，因為他是從男人身上取出來的。因此，人要離開父母
，與妻子連合，二人成為一體” （創2：18-24）。

\subsubsubsection{神設立婚姻的心意}


瑪拉基先知在這段經文指出了神設立婚姻的目的。“乃是他願人得虔誠的後裔” （瑪2：15
a）。神有靈的餘力可以造多人，但祂願意約束並限制自己的自由意志和能力，只造了一個人亞當，並且祂知道亞當會失敗並且背叛神，但祂依然願意等待，給亞當及從亞當所生出的悖逆的後裔們悔改的機會。為的是給我們做出榜樣，盼望我們效法神，按照神希望的方式正確的使用祂賜給我們的自由。以至於當末後的亞當，就是基督將救贖的道路顯明以後，我們可以憑著耶穌所指明的這條救贖的道路，重新與神和好。人透過在婚姻中信守婚約，可以經歷，體會並明白神的信實。

在《何西阿書》中，神用婚姻關係比喻祂與以色列人的關係描述神的信實和人的背約。並
且透過何西阿婚姻的痛苦經歷，使何西阿先知明白神的痛心，以及神對祂的兒女永不改變的愛。使先知可以準確的向以色列人傳講悔改的信息。“耶和華對我說：「你再去愛一個
淫婦，就是她情人所愛的；好像以色列人，雖然偏向別神，喜愛葡萄餅，耶和華還是愛他
們。」” （何西阿書 3：1）。“我必聘你永遠歸我為妻，以仁義、公平、慈愛、憐憫聘你
歸我； 也以誠實聘你歸我，你就必認識我－耶和華。”（何西阿書 2：19）。


\subsubsubsection{婚姻是人類社會一切關係的根基}

\begin{table}[H]
\centering
\begin{tabular}{p{6em}|p{4em}|p{4em}|p{4em}|p{4em}}  \toprule
 關係 & 夫妻 &子女& 父母& 弟兄姐妹\\  \midrule
創造次序&第一& 第二& 第二 &第三\\  \midrule
血緣關係&沒有 &有& 有& 有\\  \midrule
自由選擇權& 可以& 不可以 &不可以 &不可以\\  \midrule
\end{tabular}
\end{table}

a) 從創造次序上
夫妻關係在一切人類的直系親屬中是第一個創造的，是其他一切直系親屬關係的根基。當亞當和夏娃的第一個孩子該隱出生後，父母關係和子女關係同時產生了。而當他們夫妻的第二個孩子賽特出生後，人類歷史上最直系的關係就是弟兄姐妹關係隨之產生。
b) 從血緣關係上
從上面的這個表格中可以看出，在人類最直系的四個關係中，只有夫妻關係沒有血緣關
係，其他的三個關係因為有血緣關係，比較容易維護一個穩固的關係。而夫妻因為沒有直
接的血緣關係，而且男女雙方在性別上存在著很大的差異，成長與不同的家庭和不同的生長環境，與其他的直系關係相比，是屬於最不穩固，也最難維持的關係。但親密的程度卻是超越與其他的三種直系親屬關係。
c) 從自由選擇權上
雖然夫妻關係是最難維持的關係，但卻是唯一的一個神給人自由意志去選擇的一種關係。其他的三個關係，都是神替人做選擇。當神給人自由意志做出選擇後，神盼望人類學會正確的做出自己的選擇，並對自己的選擇和對配偶的承諾終身負責。

\subsubsection{聖經中神對進入婚姻盟約之前的男女的教導}

\subsubsubsection{對擇偶方面的普遍教導}
在以色列人進入迦南之前，神透過摩西，切切地吩咐他們：“不可與他們（迦南人）結親
。不可將你的女兒嫁他們的兒子，也不可叫你的兒子娶他們的女兒；因為他必使你兒子轉
離不跟從主，去事奉別神，以致耶和華的怒氣向你們發作”（申7：3-4） 。然而，以色列
人進入迦南之後，很快便“娶他們的女兒為妻，將自己的女兒嫁給他們的兒子，並事奉他
們的神。以色列人行耶和華眼中看為惡的事，忘記耶和華他們的神，去事奉諸巴力和亞舍
拉，所以耶和華的怒氣向以色列人發作”（士3：6-8）。就連世界上最有智慧的“以色列王所羅門不是在這樣的事上犯罪嗎？在多國中並沒有一王象他，且蒙他上帝所愛，上帝立
他作以色列全國的王；然而連他也被外邦女子引誘犯罪”（尼13：26）。 舊約中神透過摩西五經，多次告誡以色列人進入迦南後不可與迦南人通婚，禁止以色列人娶（嫁）拜偶像的異族人，恐怕這些外邦人誘惑以色列百姓離棄真神，去事奉外邦人的假神。而以色列人被擄的歷史告訴我們，他們沒有遵守神的吩咐，導致整個以色列南北國均被外邦人所滅。今天在新約中，聖經也多次告訴我們，“信與不信不能同負一軛”（林後6：14）。在進入婚姻之前，在主內找婚姻中的另一半，應當是我們的擇偶標準，是我們應當遵守的。順命就必蒙福。

\subsubsubsection{對喪偶者再婚的教導}

聖經中指出年輕的寡婦如果信心不能持守獨身可以再嫁，但是要找主内的弟兄。“我對
著沒有嫁娶的和寡婦說，若他們常像我就好。倘若自己禁止不住，就可以嫁娶。與其慾
火攻心，倒不如嫁娶為妙”。(弗7：8，9) 。所以我願意年輕的寡婦嫁人，生養兒女，治
理家務，不給敵人辱罵的把柄(提前5：14) 。“丈夫活著的時候，妻子是被約束的；丈夫
若死了，妻子就可以自由，隨意再嫁，只是要嫁這在主裡面的人”( 弗7：39)。

\subsubsubsection{對祭司擇偶標準的教導}
舊約中的以色列祭司擔任教導百姓分辨潔淨的和不潔淨的事物，按照神的典章判斷是非
，在聖殿中服事神，因此神對他們有更高的擇偶標準。“不可娶寡婦和被休的婦人為妻
，只可娶以色列後裔中的處女，或是祭司遺留的寡婦”。對於我們今天凡是立志在神面
前委身事祂的人，所應當遵守的準則。

\subsubsubsection{對持守獨身者的教導}
“論到你們信上所提的事，我說男不近女倒好”（林前7:1）。“我對著沒有嫁娶的和寡
婦說，若他們常像我就好（林前7:8）”。“論到童身的人，我沒有主的命令，但我既蒙
主憐恤能作忠心的人，就把自己的意見告訴你們。因現今的艱難，據我看來，人不如守
素安常才好（林前7:8）”從以上幾節經文，可以看出對於沒有結過婚，或喪偶的單身者，如果能夠持守獨身，專心是否將是最蒙福的選擇。

\subsection{聖經中神對已經進入婚姻盟約中的男女的教導}
\subsubsection{神對婚姻夫妻關係方面的教導}
“你們作妻子的，當順服自己的丈夫，如同順服主。因為丈夫是妻子的頭，如同基督是
教會的頭；他又是教會全體的救主。教會怎樣順服基督，妻子也要怎樣凡事順服丈夫。
你們作丈夫的，要愛你們的妻子，正如基督愛教會，為教會捨己。 丈夫也當照樣愛妻子
，如同愛自己的身子；愛妻子便是愛自己了。從來沒有人恨惡自己的身子，總是保養顧
惜，正像基督待教會一樣，因我們是他身上的肢體。為這個緣故，人要離開父母，與妻
子連合，二人成為一體。然而，你們各人都當愛妻子，如同愛自己一樣。妻子也當敬重
他的丈夫(弗5：22-33) ”。
從這節經文可以看出，神的心意是盼望藉著婚姻，丈夫與妻子可以合一，在婚姻中彰顯
神的榮耀。妻子學會順服丈夫，丈夫學會用基督捨己的愛去愛自己的妻子。從美好守約
的婚姻中，經歷基督對教會捨己的愛，教會對基督的順服。


\subsubsection{神對休妻方面的教導}
瑪拉基書2：16a上說的：“耶和華以色列的神說：‘休妻的事是我所恨惡的’”主耶穌也
曾說：我告訴你們，凡休妻另娶的，若不是為淫亂的緣故，就是犯姦淫了；有人娶那被休
的婦人，也是犯姦淫了。 （太5：32）儘管如此，神瞭解離婚是會發生的。在舊約中，神制定了律法來保護婦女的權利。“人若娶妻以後，見他有什麼不合理的事，不喜悅他，就可以寫休書交在他手中，打發他離開夫家。”（申24：1） 主耶穌指出這些律法是因為人心堅硬而制定，並非神想要的本意，“摩西因為你們的心硬，所以許你們休妻，但起初並不是這樣。”（太19：8）。聖經中雖然摩西允許男人休妻子，神的本意是對被棄絕的妻子的一個保護和關愛。但神對休妻的男人非常恨惡。

\subsubsection{神對離婚方面的教導}
依照聖經，神所設計的婚姻是終身的承諾。“既然如此，夫妻不再是兩個人，乃是一體的
了。所以，神配合的，人不可分開”（太19：6）。
“妻子不可離開丈夫。若是離開了，不可再嫁。或是仍同丈夫和好。丈夫也不可離棄妻子。倘若
某弟兄有不信的妻子，妻子也情願和他同住，他就不要離棄妻子。妻子有不信的丈夫，丈
夫也情願和她同住，她就不要離棄丈夫。因為不信的丈夫就因著妻子成了聖潔，並且不信
的妻子就因著丈夫成了聖潔”（林前7：10-14）。
“然而爲了保持和睦，“倘若那不信的人要離去，就由他離去吧！無論是弟兄，是姐妹，
遇著這樣的事都不必拘束。神召我們原是要我們和睦。（林前7：15）” 。
“凡休妻另娶的，若不是為淫亂的緣故，就是犯姦淫了”（太5：32）。
從這幾節經文中可以看出離婚不是神的心意。但有兩種情況，神的兒女可以離婚：1）如
果婚姻中的一方犯下了淫亂的罪，另外的一方可以提出離婚；2）如果不信主的一方要提
出離婚，為了保持和睦，信主的一方可以任憑對方離婚。

\subsubsection{神對虐妻有什麼教導} 
耶和华以色列的神说：“休妻的事和以强暴待妻的人，都是我所恨恶的！所以当谨守你们
的心，不可行诡诈，「这是万军耶和华说的」”（玛2：16） 。神及其痛恨在婚姻中實行
家庭暴力的事情，尤其是以何種方式虐待妻子的丈夫。

\subsection{瑪拉基時代以色列人婚姻的主要問題分析}

\subsubsection{瑪拉基時代以色列人婚姻中的主要問題}
瑪拉基書的寫作時間不詳。很有可能是與尼希米同時代或尼希米工作的晚期。從這段經文記錄中可以看出，被擄歸回的百姓背離了他們與神所立的約，在婚姻生活中表現在如下幾個方面：
1. 猶大人在結婚前選擇結婚的對象上，褻瀆耶和華所喜愛的聖潔，娶事奉外邦神的女子為妻，違背了摩西五經中的關於擇偶的教導。
2. 猶大人在結婚後不遵守上帝的約，用錯誤的態度處理婚姻中的問題或矛盾。以詭詐待妻，離棄幼年所娶的妻並以強暴待妻。 


\begin{table}[H]
\centering
\begin{tabular}{p{4em}|p{10em}|p{20em}|p{14em}}  \toprule
&猶大百姓的罪&神的態度&神的心意\\  \midrule
擇偶時&娶事奉外邦神的女子為妻&就是獻供物給萬軍之耶和華，耶和華也必從雅各的帳棚中剪除他。&持守信仰的聖潔\\  \midrule
結婚后&休妻的事和以強暴待妻&耶和華不再看顧那供物，也不樂意從你們手中收納&願人得虔誠的後裔\\  \midrule
\end{tabular}
\end{table}


\subsubsection{當時以色列人錯誤婚姻觀點的根源分析}
從聖經的記載中可以看出，被擄歸回的百姓犯下這樣的罪有兩方面的原因。
\subsubsubsection{歷史上的原因}
先前的君王，娶外邦女子為妻，樹立了負面的影響，百姓沒有接受教訓，反而繼續效法
。所羅門王在法老的女兒之外，又寵愛許多外邦女子，就是摩押女子、亞捫女子、以東
女子、西頓女子、赫人女子。所羅門有一千的妃嬪。所羅門年老的時候，他的妃嬪誘惑
他的心去隨從別神，不效法他父親大衛誠誠實實地順服耶和華。因為所羅門隨從西頓人
的女神亞斯她錄和亞捫人可憎的神米勒公，並行耶和華眼中看為惡的事。所羅門為摩押
可憎的神基抹和亞捫人可憎的神摩洛，在耶路撒冷對面的山上建築邱壇。他為那些向自
己的神燒香獻祭的外邦女子，就是他娶來的妃嬪也是這樣行。耶和華向所羅門發怒，因
為他的心偏離向他兩次顯現的耶和華，以色列的上帝。耶和華曾吩咐他不可隨從別神，
他卻沒有遵守耶和華所吩咐的。所以耶和華對他說：你既行了這事，不遵守我所吩咐你
守的約和律例，我必將你的國奪回，賜給你的臣子。然而，因你父親大衛的緣故，我不
在你活著的日子行這事，必從你兒子的手中將國奪回。（王上11：1-12）另外，以色列
和猶大家的君王亞哈娶了西頓人耶洗別，猶大王約蘭娶了耶洗別的女兒雅他利亞。導致
以色列和猶大的百姓開始拜巴力和亞舍拉。尤其是猶大家的王室險些被雅他利亞滅絕。

\subsubsubsection{同時代的首領和祭司及利未人對百姓的影響}
從以斯拉記中可以看到，被擄歸回的猶大首領，祭司並利未人公然帶頭違背神的律法，
娶外邦女子爲妻，“以色列民和祭司並利未人，沒有離絕迦南人、赫人、比利洗人、耶
布斯人、亞捫人、摩押人、埃及人、亞摩利人，仍效法這些國的民，行可憎的事。因他
們為自己和兒子娶了這些外邦女子為妻，以致聖潔的種類和這些國的民混雜；而且首領
和官長在這事上為罪魁。”(拉9:1-2）以斯拉帶領百姓認罪悔改，並離棄所娶的外邦女
子。然而，過了几十年，尼希米回去時，歸回的祭司中仍然有與外邦人通婚的情況。“
蒙派管理我們神殿中庫房的祭司以利亞實與多比雅結親”（尼13：4），這多比雅是亞
捫人（尼4：3），多次用各種方式破壞尼希米所領導的城墻重建工作，是猶大人的仇敵
領袖。”大祭司以利亞實的孫子、耶何耶大的一個兒子是和倫人參巴拉的女婿” （尼
13：28）。“那些日子，我也見猶大人娶了亞實突、亞捫、摩押的女子為妻”（尼
13：24）。於是尼希米再次帶領猶大百姓離棄與外邦通婚的罪。


\section{被掳归回的历史}


\subsection{巴比倫的历史}
\begin{table}[!hbp]
\centering 
\begin{tabular}{|p{3.cm}|p{1.5cm}|p{2.5cm}|p{6.2cm}|p{2.5cm}|}
\hline
\hline
巴比倫王 & 做王时间 &参考出处 & 歷史事件 &备注 \\\hline
那波帕拉沙 Nabopolassar& 626(630?)-605& \cite{Nabopolassar} & 本是一个迦勒底王子，和米底亞國王庫阿克斯列斯（Huvakhstra）聯手，在 612BC 圍攻尼尼微，609BC亞述滅亡&“那波帕拉薩爾” “拿布普拉撒” \\\hline
&609 &王下23：29，代下35：24 & 約西亞王被埃及王尼哥在米吉多所杀& \\\hline
&609 &王下23：31-34，代下36：1-4 & 约哈斯被法老擄到埃及& \\\hline
&609 & 王下23：34,代下36：4& 以利雅敬作猶大和耶路撒冷的王，改名叫約雅敬&約哈斯的哥哥 \\\hline
&606& 但1：1-2&尼布甲尼撒將耶路撒冷圍困,把神殿中器皿的幾分帶到他神的庫中&\\\hline
&606 & 但1：3-6&但以理,哈拿尼雅、米沙利、亞撒利雅被掳到巴比伦&\\\hline
\textbf{尼布甲尼撒 Nebuchadnezzar}& 605-561&  & & 那波帕拉沙兒子 \\\hline
%&605 & &耶25：1& 耶利米在约雅敬第四年预言犹大要被掳至巴比伦七十年& \\\hline
%&604 & &但2:1& 但以理给尼布甲尼撒解梦& \\\hline
%&?  & &但3:1& 尼布甲尼撒造金像& \\\hline
%& ? & &但4：1& 但以理给尼布甲尼撒解第二个梦& \\\hline
%&  ?& &但4：29, 33& 尼布甲尼撒被趕出離開世人& \\\hline
& 598 & 代下36：6 & 约雅敬（609-598）被掳至巴比伦 & \textcolor[rgb]{0,.1,1}{\textbf{第一次被掳}}\\\hline
& 597 &  代下36：10& 尼布甲尼撒差遣人將約雅斤（以西结）和耶和華殿裡寶貴器皿帶到巴比倫 &\textcolor[rgb]{0,.1,1}{\textbf{第二次被掳}} \\\hline
& 597 & 代下36：10& 西底家作猶大和耶路撒冷的王 & \\\hline
&586 & 代下36：11& 西底家（Zedekiah 597-586BC）被掳巴比伦&\textcolor[rgb]{0,.1,1}{\textbf{第三次被掳}}\\\hline
\textbf{以未米罗达 Evilmerodach}& 562-560 && & 尼布甲尼撒之子\\\hline
	&562	&	王下25：27&	把约雅斤（Jehoiachin）释放& \\\hline

涅里格利沙爾 Neriglissar & 560-556 & &在前560年殺害了以未米羅達而登位 & 尼布甲尼撒的女婿\cite{Neriglissar} “尼甲沙利薛”\cite{尼甲沙利薛}\\\hline
拉巴施馬爾杜克 Labashi-Marduk&  556& & 巴比倫涅里格利沙爾之子，幼年繼位，在位不足一年便被殺害\cite{Labashi-Marduk}&\\\hline
拿波尼度 Nabonidus&	555-539&	&那波尼德大約在前549年離開巴比倫，任命其子伯沙撒為共同攝政\cite{Nabonidus}& 拿波尼度娶了尼布甲尼撒女儿李道葵斯\cite{尼甲沙利薛}\\\hline
\textbf{伯沙撒	Belshazzar}&553-539&	但5：11，18& & 伯沙撒的母亲相传是尼布甲尼撒的女儿\cite{尼甲沙利薛}\\\hline	

	%		耶51：11	巴比伦帝国要被瑪玳所灭“耶和华定意攻击巴比伦，将她毁灭，所以激动了玛代（米底亚）君王的心，因这是耶和华报仇，就是为自己的殿报仇。”
&553&但7& 伯沙撒元年,但以理見四兽的異象& \\\hline	
	&551\cite{伯沙撒}&但8&伯沙撒第三年,但以理見末后四国異象& \\\hline
		&539  &	 	 但5& 伯沙撒被殺當夜,但以理预言亡国& \\\hline	
\end{tabular}
\caption{被掳归回的历史 (巴比倫) 1}
\end{table}

 
\subsection{玛代/波斯的历史}
\begin{table}[!hbp]
\centering 
\begin{tabular}{|p{2.5cm}|p{1.7cm}|p{1.5cm}|p{6.cm}|p{2.8cm}|}
\hline
\hline
玛代/波斯王 & 做王时间 &参考出处 & 歷史事件 &备注 \\\hline
\textbf{大利烏（玛代）}&	539-538	&	 	 &大利乌成了古列分封的巴比伦王\cite{伯沙撒}	&   \\\hline
&	元年538	&	但9	 &   但以理得知耶路撒冷荒涼七十年。便禁食，向神祈禱得知七十個七的启示&   \\\hline
&	元年538	&	但11：1	 &   我曾起來扶助米迦勒，使他堅強。&   \\\hline
 &	 	&	但6：28	 &   但以理当大利乌王在位的时候和波斯王塞鲁士在位的时候，大享亨通。	&   \\\hline
&	 	&	但6	 &   但以理被投入狮子坑	&   \\\hline
\textbf{古列 Cyrus    （波斯）}&	559-529	&代下36: 22 拉1:1 &（爷爷波斯王）鐵伊司佩斯（父亲）岡比西斯一世，（母亲）米底亞王的女兒 & 又名“塞鲁士”， “居魯士二世” \\\hline	
&538（7）\cite{居鲁士}&拉1:1 & 波斯王古列元年下诏，重建圣殿&   \\\hline	
&536&但10-12 & 第三年但以理得末后大争战的启示&   \\\hline	
&536& &所罗巴伯回归耶路撒冷& \textcolor[rgb]{0,.1,1}{\textbf{第一次歸回}} \\\hline	
冈比西斯二世 Cambyses&	529-522	&  &居鲁士大帝的儿子，冈比西斯二世将兄弟巴尔迪亚杀死继王位  &   \\\hline	
高墨达\cite{高墨达}&	522（3或8个月）&& 伪称自己是巴尔迪亚，夺取了王位。玛代人斯默迪斯，犹太人称为亚达薛西 &斯梅尔迪士，  亚达薛西    \\\hline	
\textbf{大流士}& 522-486&  &居鲁士大帝的女婿& 大利乌，Darius I\\\hline	
& 520&&大利乌王下旨继续建殿的工作& \\\hline
&516&拉6：15&殿在大流士六年亚达月初三完成& \\\hline
薛西一世  & 486 &&&Xerxes I \\\hline
\textbf{亞哈隨魯}&	485-465	&	拉四：6 \cite{亞哈隨魯} &亚哈随鲁登基，敌人上本控告犹大和耶路撒冷的居民。大流士一世与居鲁士大帝之女阿托莎的儿子&“薛西斯一世”  “澤克西斯一世” “泽尔士一世” \\\hline	
&480　&但11:2 &波斯還有三王興起，第四王攻擊希臘&\\\hline
&479(8)　& & 以斯帖成为王后&\\\hline
\textbf{亚达薛西一世” Artaxerxes I} &	465-425&拉四：7-16&澤克西斯一世之子，瓦实提的儿子&阿尔塔薛西斯一世\\\hline
&458&&以斯拉回归耶路撒冷	&\textcolor[rgb]{0,.1,1}{\textbf{第二次歸回}}\\\hline
&444&&尼希米到达耶路撒冷，城墙完工	&\textcolor[rgb]{0,.1,1}{\textbf{第三次歸回}}\\\hline
薛西斯二世	&425&&亚达薛西一世嫡子，在位45天被暗杀&\\\hline	
塞基狄亚努斯 	&423&&阿尔塔薛西斯一世的庶子&\\\hline	
大流士二世	&423-404&&夺取异母兄弟塞基狄亚努斯王位&\\\hline
&420&&尼希米再到耶路撒冷	&\\\hline
亚达薛西二世&404-359&&大流士二世儿子&Artaxerxes II\\\hline
亚达薛西三世&358-338&&&Artaxerxes III\\\hline
阿尔塞斯&338-336&&亚达薛西三世的幼子& \\\hline
大流士三世&336-330&&亚达薛西二世的侄孫，被貝蘇斯所杀&Darius III\\\hline
貝蘇斯&330-329&&亞歷山在巴克特里亞處死了貝蘇斯& \\\hline
%拉六：3，14  塞鲁士王元年，他降旨论到耶路撒冷上帝的殿，要建造这殿为献祭之处，坚立殿的根基。殿高六十肘，宽六十肘。。。犹大长老因先知哈该和易多的孙子撒迦利亚所说劝勉的话，就建造这殿，凡事亨通。他们遵着以色列上帝的命令和波斯王塞鲁士、大利乌（或大流士一世 Darius I or Darius the Great，公元前522-前486年)、亚达薛西（Artaxerxes I，公元前465-前425年)的旨意，建造完毕。								
%在先知书上，我们也可以看到他的名字：									
%但一：21  到塞鲁士王元年，但以理还在。									
								
%但十：1  波斯王塞鲁士第三年，有事显给称为伯提沙撒的但以理。这事是真的，是指着大争战。但以理通达这事，明白这异象。														
%但最令人惊讶的是，先知以赛亚竟然预言塞鲁士在历史舞台上的出现，这是不可%思议的事，因为以赛亚是在“乌西雅（Uzziah 776/775-736/735BC）、约坦（Jotham 750-735/730BC）、亚哈斯（Ahaz 735/734 or 731/730-715BC）、希西家（Hezekiah 715-687/686BC）作犹大王的时候得默示”（赛一：1）这是在塞鲁士（Cyrus II 559-529BC）出现之前约两百年。有可能吗？不信派的人当然说不可能，特别是名字都出现在《以赛亚书》三十九章之后：									
%赛四十四：28  论塞鲁士说：“他是我的牧人，必成就我所喜悦的，必下令建造耶路撒冷，发命立稳圣殿的根基。”									
%赛四十五：1  我耶和华所膏的塞鲁士，我搀扶他的右手，使列国降伏在他面前。我也要放松列王的腰带，使城门在他面前敞开，不得关闭。我对他如此说：。。						赛四十五：13  我凭公义兴起塞鲁士，又要修直他一切道路。他必建造我的城，释放我被掳的民，不是为工价，也不是为赏赐。这是万军之耶和华说的。						
%还有赛四十一：2-3，25-26，四十六：11 都暗示着“塞鲁士”的到来。								
\end{tabular}
\caption{被掳归回的历史(玛代/波斯) 2}
\end{table}

\subsection{希臘的历史}
\begin{table}[!hbp]
\centering 
\begin{tabular}{|p{2.7cm}|p{1.4cm}|p{11.2cm}|}
\hline
\hline
希臘王 & 时间 & 歷史事件 \\\hline
亚历山大大帝&336-323  & 繼其父腓力二世王位成為馬其頓國王。公元前334年，征服波斯帝國。 \\\hline
巴比倫分封協議&323&  \\\hline
%希臘叛亂 &323-322 &   \\\hline
%繼業者之戰&320-301&     \\\hline
%&320   &第一次繼業者之戰 \\\hline
%&319-315   &第二次繼業者之戰 \\\hline
%&314-311   &第三次繼業者之戰 \\\hline
%&308-301  &第四次繼業者之戰 \\\hline
%馬其頓的鬥爭&298-285  &  \\\hline
%利西馬科斯與塞琉古之爭&285-281  &  \\\hline
%凯尔特人入%侵&285-275 &  \\\hline
%皮洛士最後一搏&274-272 &  \\\hline
%敘利亞戰爭&274-101 &  \\\hline
 %&脫拉西及小亞西亞（賴西馬克Lysimachus）
%&埃及和巴勒斯坦的多利買一世、&  \\\hline
%&巴比倫和敘利亞的西流基&  \\\hline
%&小亞細亞的呂西馬古及馬其頓&  \\\hline
安提柯王朝&323-168&馬其頓和希臘, Cassander \\\hline
托勒密王朝 &323-30&埃及和巴勒斯坦地區（Ptolemy多利買） \\\hline
塞琉古王朝 &323-64&美索不達米亞、波斯、小亞細亞（Seleucus I Nicator，西流基） \\\hline
&323-303&利西馬邱斯(Lysimachus) 控制小亞細亞及色雷斯，20年後，四國減為三國，塞流卡斯殺了利西馬邱斯，併吞小亞細亞\\\hline
 \end{tabular}
\caption{希臘王\cite{繼業者}}
\end{table}

 

 \section{以色列月份}
\begin{table}[!hbp]
\centering 
\begin{tabular}{|p{1.2cm}|p{7cm}|p{3.6cm}|p{2.8cm}|}
\hline
\hline
%http://www.pcchong.com/Esther/Esther6.htm
月份　&以色列月份&聖經章節&陽曆時間\\\hline
正月 &尼散月 Nisan， 起初叫亚笔月 Abib&出十三：4， 出十二：2&阳历的三/四月\\\hline
二月 & 以珥月（或西弗月 Ziv，起初叫基乌月 Iyyar）&王上六：1&阳历四/五月\\\hline
三月 &西弯月（Sivan）&&阳历五/六月\\\hline
四月 & 塔模斯月（Tammuz）&&阳历六/七月\\\hline
五月& 埃波月（Av）&&阳历七/八月\\\hline
六月& 以禄月（Elul）&&阳历八/九月\\\hline
七月& Tishri（起初叫以他念月 Ethanim）&王上八：2&阳历九/十月\\\hline
八月 &- Marchesvan （起初叫布勒月 Bul）&王上六：38&阳历十/十一月\\\hline
九月& 基斯流月（Kislev）&&阳历十一/十二月\\\hline
十月 & 提别月（Tevet）&&阳历十二/一月\\\hline
十一月 & 示巴特月（Shvat）&&阳历一/二月\\\hline
十二月&亚达月（Adar）&&阳历二/三月\\\hline
\end{tabular}
\caption{以色列月份}
\end{table}


 \section{被掳归回与先知书}
\subsection{被掳归回与但以理书}
\begin{table}[!hbp]
\centering 
\begin{tabular}{|p{2.8cm}|p{1.5cm}|p{2.5cm}|p{8.5cm}| }
\hline
\hline
巴比倫王 &  时间 &参考出处 & 歷史事件   \\\hline
那波帕拉沙 & 626(630?)-605& \cite{Nabopolassar}   &Nabopolassar, “那波帕拉薩爾” “拿布普拉撒” \\\hline

&606 & 但1：6,7&但以理,哈拿尼雅、米沙利、亞撒利雅被掳到巴比伦 \\\hline
\textbf{尼布甲尼撒}& 605-561&  &Nebuchadnezzar, 那波帕拉沙兒子 \\\hline
&604 & 但1：18-20&但以理,哈拿尼雅、米沙利、亞撒利雅被留在王面前  \\\hline
%&605 & &耶25：1& 耶利米在约雅敬第四年预言犹大要被掳至巴比伦七十年& \\\hline
&604 & 但2& \textcolor[rgb]{0,.1,1}{\textbf{但以理给王解梦}},王派他管理巴比倫全省，立他為總理，掌管巴比倫的一切哲士。派沙得拉、米煞、亞伯尼歌管理巴比倫省的事務，但以理常在朝中侍立。  \\\hline
&?  &  但3& 尼布甲尼撒造金像並降旨，任何人不得謗讟神，在巴比倫省，高陞了沙得拉、米煞、亞伯尼歌。  \\\hline
& 597 &  代下36：10& 尼布甲尼撒差遣人將約雅斤（以西结）和耶和華殿裡寶貴器皿帶到巴比倫 ,第二次被掳 \\\hline
%& 597 & 代下36：10& 西底家作猶大和耶路撒冷的王 & \\\hline
&586 & 代下36：11& 西底家（Zedekiah 597-586BC）被掳巴比伦\\\hline
& ? &但4：1-27& 但以理给尼布甲尼撒解第二个梦  \\\hline
&12月後 &  但4：29, 33& 尼布甲尼撒被趕出離開世人七年  \\\hline
&過了七年 &  但4：34-37& 聰明復歸於尼布甲尼撒 \\\hline
%& 598 & 代下36：6 & 约雅敬（609-598）被掳至巴比伦 & 第一次被掳\\\hline

\textbf{以未米罗达 }& 562-560 &  &Evilmerodach, 尼布甲尼撒之子，释放约雅斤\\\hline
	%&562	&	王下25：27&	把约雅斤（Jehoiachin）释放& \\\hline
涅里格利沙爾  & 560-556 & &Neriglissar, 尼布甲尼撒的女婿\cite{Neriglissar} “尼甲沙利薛”\cite{尼甲沙利薛}\\\hline
拉巴施馬爾杜克 &  556&  &Labashi-Marduk,涅里格利沙爾之子，幼年繼位\cite{Labashi-Marduk} \\\hline
拿波尼度 &	555-539&	&Nabonidus,\cite{Nabonidus}, 娶了尼布甲尼撒女儿\cite{尼甲沙利薛}\\\hline
\textbf{伯沙撒	}&553-539&	但5：11，18&Belshazzar,  母亲相传是尼布甲尼撒的女儿\cite{尼甲沙利薛}\\\hline	
	%		耶51：11	巴比伦帝国要被瑪玳所灭“耶和华定意攻击巴比伦，将她毁灭，所以激动了玛代（米底亚）君王的心，因这是耶和华报仇，就是为自己的殿报仇。”
&553&但7& \textcolor[rgb]{0,.1,1}{\textbf{伯沙撒元年,在床上做夢，見了腦中的四兽的異象}}  \\\hline	
	&551\cite{伯沙撒}&但8&\textcolor[rgb]{0,.1,1}{\textbf{伯沙撒第三年,但以理在以攔省書珊城中；見公綿羊公山羊異象又如在烏萊河邊。 加百列解釋異象}}  \\\hline
		&539  &	 	 但5& 伯沙撒被殺當夜,但以理预言亡国  \\\hline			
\textbf{大利烏（玛代）}&	539-538	&	 	 &大利乌成了古列分封的巴比伦王\cite{伯沙撒}	    \\\hline
&	元年538	&	但9	 &   \textcolor[rgb]{0,.1,1}{\textbf{但以理得知耶路撒冷荒涼七十年。便禁食，向神祈禱.加百列奉命傳講七十個七的启示 }}   \\\hline
&	元年538	&	但11：1	 &   我曾起來扶助米迦勒，使他堅強。   \\\hline
% &	 	&	但6：28	 &   但以理当大利乌王在位的时候和波斯王塞鲁士在位的时候，大享亨通。	    \\\hline
&	 	&	但6	 &   但以理被大利乌王投入狮子坑	    \\\hline
\textbf{古列     （波斯）}&	559-529	& 拉1:1 &Cyrus, 又名“塞鲁士”， “居魯士二世” \\\hline	
&538/7 \cite{居鲁士}&拉1:1 & 波斯王古列元年下诏，重建圣殿    \\\hline	
&536&但10-12 & \textcolor[rgb]{0.0,.1,1}{\textbf{第三年但以理在希底結大河邊,三位天使向他傳達末后大争战的启示} }   \\\hline			
\end{tabular}
\caption{但以理}
\end{table}

\subsection{被掳归回与耶利米书}
\begin{table}[!hbp]
\centering 
\begin{tabular}{|p{1.8cm}|p{2cm}|p{1.5cm}|p{1.7cm}|p{8.5cm}| }
\hline
\hline
巴比倫王 &以色列王  & 时间 &参考出处 & 歷史事件   \\\hline
坎达拉努& &647-627 & &巴比伦第九王朝末位国王,他由亚述拥立。 \\\hline
&約西亞在位十三年&627   &耶1：2&耶和華的話臨到耶利米。''立你在列邦列國之上，為要施行拔出、拆毀、毀壞、傾覆，又要建立、栽植。''  \\\hline
那波帕拉沙 & &626-605& \cite{Nabopolassar}   &他在巴比伦国王坎达拉努死后夺取王位,于公元前612年占领亚述都城尼尼微。Nabopolassar \\\hline
 & 約西亞王在位的時候 &  &耶3:6  &  責備背道的以色列因行淫被棄絕,他奸詐的妹妹猶大不懼怕，也去行淫。 \\\hline
 
  ?& 約雅敬&  &耶22:11  &約哈斯必不得再回到這裡來，卻要死在被擄去的地方，必不得再見這地 \\\hline
  ?& 約雅敬&  &耶22:18&  約雅敬被埋葬，好像埋驢一樣，要拉出去扔在耶路撒冷的城門之外。 \\\hline
    ?& 約雅敬&  &耶22:24  &  耶哥尼雅必將你交給尼布甲尼撒和迦勒底人的手中。也必將你和生你的母親趕到別國，你們必死在那裡，不得歸回。\\\hline
那波帕拉沙 &約雅敬第一年 &609&耶26:1  &你們若不聽從我，我就必使這殿如示羅，使這城為地上萬國所咒詛的。               \\\hline  
    尼布甲尼撒 & 約雅敬第四年& 605 &耶25:1  &這全地必然荒涼，令人驚駭。這些國民要服事巴比倫王七十年。滿了以後，我必刑罰巴比倫王和那國民，並迦勒底人之地，因他們的罪孽使那地永遠荒涼。 \\\hline
       尼布甲尼撒 & 約雅敬第四年& 605 &耶36:1-8  &耶利米召了尼利亞的兒子巴錄來；將耶和華對耶利米所說的一切話寫在書卷上。沙番的孫子、基瑪利雅的兒子\underline{米該亞} 向眾首領，文士\underline{以利沙瑪} 、示瑪雅的兒子\underline{第萊雅} 、亞革波的兒子\underline{以利拿單} 、沙番的兒子\underline{基瑪利雅} 、哈拿尼雅的兒子\underline{西底家} ，和其餘的首領講耶利米的話。眾首領就打發古示的曾孫、示利米雅的孫子、尼探雅的兒子\underline{猶底}到巴錄那裡拿書卷來到他們那裡。猶底將這一切話說給王聽。王就用文士的刀將書卷割破，扔在火盆中燒盡了。 以利拿單和第萊雅，並基瑪利雅懇求王不要燒這書卷。王就吩咐哈米勒的兒子\underline{耶拉篾}和亞斯列的兒子\underline{西萊雅}，並亞伯疊的兒子\underline{示利米雅}，去捉拿巴錄和耶利米。耶利米寫了約雅敬所燒前卷上的一切話交給巴錄。 \\\hline
  
  
   尼布甲尼撒 & 約雅敬第四年& 605 &耶45   &神應許拯救尼利亞的兒子巴錄\\\hline
   
         尼布甲尼撒 & 約雅敬第五年九月& 604 &耶36   &巴錄就在沙番的兒子文士\underline{基瑪利雅}的屋內，念書上耶利米的話給眾民聽。\\\hline
      
       
  尼布甲尼撒& 約雅敬&  &耶35  &耶利米對利甲族的人說：利甲的兒子約拿達必永不缺人侍立在我面前。  \\\hline    
   
\end{tabular}
\caption{耶利米 1}
\end{table}



\begin{table}[!hbp]
\centering 
\begin{tabular}{|p{1.8cm}|p{2cm}|p{1.5cm}|p{1.7cm}|p{8.5cm}| }
\hline
\hline
巴比倫王 &以色列王  & 时间 &参考出处 & 歷史事件   \\\hline
尼布甲尼撒& 西底家王&598  &耶24:1 &一筐是極好的無花果；一筐是極壞的無花果，壞得不可吃  \\\hline 
 尼布甲尼撒 &西底家第一年 &598&耶27:1  &做繩索與軛藉那些來到耶路撒冷見猶大王西底家的使臣之手，把繩索與軛送到以東王、摩押王、亞捫王、推羅王、西頓王那裡，列國都必服事巴比倫王尼布甲尼撒              \\\hline    
 

 
 尼布甲尼撒 &西底家第四年五月 &594&耶28:1  & 耶利米預言耶和華將鐵軛加在列國的頸項上，使他們服事巴比倫王尼布甲尼撒   \\\hline 
 尼布甲尼撒 &西底家第四年七月 &594&耶28:17  & 基遍人押朔的兒子先知\underline{哈拿尼雅}將先知耶利米頸項上的軛取下來，折斷了。 當年七月間就死了   \\\hline  
  尼布甲尼撒 &西底家第四年  &594&耶51  & 西底家上巴比倫去，瑪西雅的孫子、尼利亞的兒子西萊雅與王同去,耶利米讓西萊雅到了巴比倫務要念這書上一切要臨到巴比倫的災禍  \\\hline 
 尼布甲尼撒& 西底家王&  &耶21:1  & 耶和華的話臨到耶利米。西底家王打發瑪基雅的兒子\underline{巴施戶珥}和瑪西雅的兒子祭司\underline{西番雅}去見耶利米\\\hline
尼布甲尼撒& 西底家王&  &耶29  & 耶利米藉沙番的兒子\underline{以利亞薩}和希勒家的兒子\underline{基瑪利}的手寄信與被擄的人:''為巴比倫所定的七十年滿了以後，使你們仍回此地。 論到哥賴雅的兒子\underline{亞哈}，並瑪西雅的兒子\underline{西底家}，他們是託我名向你們說假預言的，我必將他們交在巴比倫王尼布甲尼撒的手中殺害他們。我必刑罰被擄的尼希蘭人\underline{示瑪雅}和他的後裔''\\\hline

尼布甲尼撒& 西底家王&   &耶37  &示利米雅的兒子\underline{猶甲}和祭司瑪西雅的兒子\underline{西番雅}去見先知耶利米,耶和華此說:''法老的軍隊必回埃及本國去。迦勒底人必再來攻取這城，用火焚燒'' \\\hline

尼布甲尼撒& 西底家王&   &耶37  &哈拿尼亞的孫子、示利米雅的兒子守門官\underline{伊利雅拿}住先知耶利米將他囚在文士\underline{約拿單}的房屋中 \\\hline

尼布甲尼撒& 西底家王&   &耶38  &瑪坦的兒子\underline{示法提雅}、巴施戶珥的兒子\underline{基大利}、示利米雅的兒子\underline{猶甲}、瑪基雅的兒子\underline{巴示戶珥}拿住耶利米，下在哈米勒的兒子\underline{瑪基雅}的牢獄裡\\\hline

尼布甲尼撒& 西底家王&   &耶38  &古實人\underline{以伯米勒}用繩子將耶利米從牢獄裡拉上來。耶利米仍在護衛兵的院中。\\\hline

尼布甲尼撒& 西底家王&   &耶38  &耶利米勸西底家出去歸降巴比倫王的首領\\\hline
尼布甲尼撒 &西底家  & &耶39  & 耶利米還囚在護衛兵院中的時候，耶和華應許拯救古實人以伯米勒 \\\hline 

 尼布甲尼撒 &西底家九年十月初十日 &588&耶39,52:4  & 巴比倫王尼布甲尼撒率領全軍來圍困耶路撒冷。\\\hline 
 

\end{tabular}
\caption{耶利米 2}
\end{table}


\begin{table}[!hbp]
\centering 
\begin{tabular}{|p{1.8cm}|p{2cm}|p{1.5cm}|p{1.7cm}|p{8.5cm}| }
\hline
\hline
巴比倫王 &以色列王  & 时间 &参考出处 & 歷史事件   \\\hline


 尼布甲尼撒 &西底家第十年 &587&耶32  & 耶利米囚在護衛兵的院內, 向他叔叔的兒子\underline{哈拿篾}買了亞拿突的那塊地，平了十七舍客勒銀子給他。  \\\hline 
 尼布甲尼撒 &西底家第十年 &587&耶33  & 耶利米還囚在護衛兵的院內，耶和華的話第二次臨到他說在那日子猶大必得救，耶路撒冷必安然居住,大衛必永不斷人坐在以色列家的寶座上    \\\hline  
 尼布甲尼撒& 西底家王&587  &耶34  &巴比倫王的軍隊正攻打耶路撒冷，又攻打猶大的拉吉和亞西加。耶利米告訴西底家，耶和華要將這城交付巴比倫王的手。西底家王與耶路撒冷的眾民立約，向他們宣告自由，卻又反悔. 耶和華說：我必吩咐巴比倫王軍隊回到這城，攻打這城，將城攻取，用火焚燒。 \\\hline
 尼布甲尼撒 &西底家第十一年四月初九日 &586&耶39, 52:6,12 &巴比倫王的首領尼甲沙利薛、三甲尼波、撒西金、拉撒力、尼甲沙利薛拉墨，並巴比倫王其餘的一切首領都來坐在中門。城被攻破。 \\\hline  

尼布甲尼撒 &西底家第十一年  &586&耶39:14,40:1  &護衛長尼布撒拉旦和尼布沙斯班拉撒力、尼甲沙利薛拉墨，並巴比倫王的一切官長，將耶利米從護衛兵院中提出來，將他從拉瑪釋放以後,交與沙番的孫子亞希甘的兒子基大利 \\\hline 

尼布甲尼撒 &西底家第十一年  &586&耶40  &軍長尼探雅的兒子以實瑪利，加利亞的兩個兒子約哈難和約拿單，單戶篾的兒子西萊雅，並尼陀法人以斐的眾子，瑪迦人的兒子耶撒尼亞和屬他們的人，都到米斯巴見基大利。 \\\hline

尼布甲尼撒 &西底家第十一年七月間  &586&耶41  &王的大臣宗室以利沙瑪的孫子、尼探雅的兒子以實瑪利帶著十個人，來到米斯巴用刀殺了沙番的孫子亞希甘的兒子基大利\\\hline

尼布甲尼撒 &西底家第十一年七月間  &586&耶41  &第二天，有八十人從示劍和示羅，並撒瑪利亞來，要奉到耶和華的殿。尼探雅的兒子以實瑪利出米斯巴迎接他們將他們殺了，拋在坑中。這坑是從前亞撒王因怕以色列王巴沙所挖的。 \\\hline
尼布甲尼撒 &西底家第十一年七月間  &586&耶41  &尼探雅的兒子以實瑪利和八個人從基遍的大水旁脫離約哈難的手，逃往亞捫人那裡去了。加利亞的兒子約哈難和同著他的眾軍長將，從米斯巴將剩下的一切百姓帶到靠近伯利恆的金罕寓住下，要進入埃及去\\\hline

尼布甲尼撒 &西底家第十一年  &586&耶42  &何沙雅的兒子亞撒利雅和加利亞的兒子約哈難不聽從耶和華的話住在猶大地,將所剩下的猶大人先知耶利米，以及尼利亞的兒子巴錄，都帶入埃及地，到了答比匿\\\hline



尼布甲尼撒 &西底家第十一年  &586&耶44  &耶利米預言 一切住在埃及地的猶大人，就是住在密奪、答比匿、挪弗、巴忒羅境內的猶大人，那定意進入埃及地、在那裡寄居的猶大人，我必使他們盡都滅絕\\\hline
\end{tabular}
\caption{耶利米 3}
\end{table}

\subsection{被掳归回与以賽亞/耶利米书}
\begin{table}[!hbp]
\centering 
\begin{tabular}{|p{3cm}|p{ 2cm}|p{2cm}|p{ 2cm}|p{2cm}|p{ 2cm}| }
\hline
\hline

國家 &以賽亞　&耶利米 \\\hline
亞述 &10：24-34，14：24-27　& \\\hline
  埃及王法老尼哥& & 46:1-26\\\hline
  雅各 & 14：1-2& 46:27-28 \\\hline
   非利士人&14：28-27 & 47\\\hline

 摩押& & 48\\\hline
 亞捫& & 49:1-6  \\\hline
  以東& & 49:7-22  \\\hline
 
  大馬色&& 49:23-27\\\hline
  基達&& 49:28-29  \\\hline 
  夏瑣& &49:30-33\\\hline
  
 以攔 & &49:34-39\\\hline
   巴比倫 &13，14：3-23& 50-51:58\\\hline 
 
\end{tabular}
\caption{耶和華論列國的話}
\end{table}


 %\begin{figure}[!htb]  
%	 \begin{center}
%%	 \includegraphics[width=3.2in]{shijian}
	% \caption{尼探雅的兒子以實瑪利}
	% \label{fig:1} 
%	 \end{center}
	% \end{figure}


\begin{thebibliography}{1}
  \bibitem{Nabopolassar} https://zh.wikipedia.org/wiki/%E9%82%A3%E6%B3%A2%E5%B8%95%E6%8B%89%E8%96%A9%E7%88%BE
  \bibitem{Neriglissar} https://zh.wikipedia.org/wiki/%E6%B6%85%E9%87%8C%E6%A0%BC%E5%88%A9%E6%B2%99%E7%88%BE

  \bibitem{Labashi-Marduk} https://zh.wikipedia.org/wiki/%E6%8B%89%E5%B7%B4%E6%96%BD%E9%A6%AC%E7%88%BE%E6%9D%9C%E5%85%8B
  
  \bibitem{Nabonidus}  https://zh.wikipedia.org/wiki/%E9%82%A3%E6%B3%A2%E5%B0%BC%E5%BE%B7
  
  \bibitem{亞哈隨魯}   https://zh.wikipedia.org/wiki/%E4%BA%9E%E5%93%88%E9%9A%A8%E9%AD%AF

  \bibitem{伯沙撒}  http://www.sdacn.org/htmlshuji/yuyan/daniel-m/9.htm

   \bibitem{尼甲沙利薛}  http://www.jiduzhijia.com/cj/Old%20Testament/27Dan/27TT00.htm
   
   \bibitem{居鲁士}  https://wapbaike.baidu.com/item/%E4%BB%A5%E6%96%AF%E6%8B%89

   \bibitem{高墨达}  https://zh.wikipedia.org/wiki/%E9%AB%98%E5%A2%A8%E8%BE%BE
   \bibitem{繼業者}  https://zh.wikipedia.org/wiki/%E7%B9%BC%E6%A5%AD%E8%80%85
  %\bibitem{notes} John W. Dower {\em Readings compiled for History
  %21.479.}  1991.

  %\bibitem{impj}  The Japan Reader {\em Imperial Japan 1800-1945} 1973:
 % Random House, N.Y.

  %\bibitem{norman} E. H. Norman {\em Japan's emergence as a modern
  %state} 1940: International Secretariat, Institute of Pacific
  %Relations.

  %\bibitem{fo} Bob Tadashi Wakabayashi {\em Anti-Foreignism and Western
 % Learning in Early-Modern Japan} 1986: Harvard University Press.


  \end{thebibliography}

